\scanpage{55}
% archive foliation: 57
$U\longrightarrow\overline H^+$ sont exactement les cartes topologiques
(i.e. \deleted{applications continues injectives,}
\inserted{plongements \unsure{topologiques} \deleted{\unsure{ouverts}}\unsure{--- ouverts et topologiques|--- ouverts et continus}}
appliquant bord sur bord) qui sont conformes dans $\operatorname{Int}U$.
Cela \deleted{\unsure{résulte du}}\inserted{équivaut au}

\uline{Lemme 1} Soit $U$ un ouvert de $\overline H^+$, alors \deleted{toute}
\deleted{application injective}
\inserted{plongement ouvert $f:U\longrightarrow\overline H$}
\deleted{fonction continue sur $V$} qui est hol.
dans $\unsure{U|V}\cap H$, est holo. (i.e. se prolonge holomorphiquement
dans $\mathbb C$, au voisinage de tt
pt de $U\cap\delta H=U\cap\mathbb R$).

C'est un cas particulier du \uline{lemme de
\unsure{Schwartz|Schwarz}}, qui dit que tt fonction \deleted{hol.}
continue complexe $f$ dans $U$, qui est
\uline{réelle} sur $U\cap\mathbb R$, est holomorphe.

Pour terminer la dem du \deleted{lemme}\inserted{th}, il faut
voir encore que la $f$ du lemme est bien
(au voisinage de tt pt $x$ de $U\cap\delta H$) une
carte holomorphe de \deleted{\unsure{$V$}} $U$ (i.e. de $H$) en $x$,
i.e. que le germe \inserted{en $x$} de fonction holomorphe
\uline{sur $\mathbb C$} défini par $f$ est inversible. Mais \unsure{ce}
germe n'est pas constant, \deleted{de la structure}\inserted{de par}\ \unsure{son} développe-
